\section{Two-Point Boundary Value Problems}

\noindent
In Section~5.3 we considered initial value problems for the
linear second order equation
\begin{equation} \label{eq:13.1.1}
P_{0}(x)y''+P_{1}(x)y'+P_{2}(x)y=F(x).
\end{equation}
Suppose $P_{0}$, $P_{1}$, $P_{2}$, and $F$ are continuous and $P_{0}$
has no zeros
on an
open interval $(a,b)$. From Theorem~\ref{thmtype:5.3.1}, if $x_{0}$ is in
$(a,b)$ and  $k_{1}$ and $k_{2}$ are arbitrary real numbers then
\eqref{eq:13.1.1} has a unique solution on $(a,b)$ such that $y(x_{0})=k_{1}$
and $y'(x_{0})=k_{2}$. Now we consider a different problem for
\eqref{eq:13.1.1}.

{\bf PROBLEM}\quad
Suppose $P_{0}$, $P_{1}$, $P_{2}$, and $F$ are  continuous and  $P_{0}$
has no zeros on a closed interval $[a,b]$. Let $\alpha$, $\beta$,
$\rho$, and $\delta$ be real numbers such that
\begin{equation} \label{eq:13.1.2}
\alpha^{2}+\beta^{2}\ne 0\text{\;    and\; \;  } \rho^{2}+\delta^{2}\ne 0,
\end{equation}
and let $k_{1}$ and $k_{2}$ be arbitrary real numbers. Find a solution of
\begin{equation} \label{eq:13.1.3}
P_{0}(x)y''+P_{1}(x)y'+P_{2}(x)y=F(x)
\end{equation}
on the closed interval $[a,b]$ such that
\begin{equation} \label{eq:13.1.4}
\alpha y(a)+\beta y'(a)=k_{1}
\end{equation}
and
\begin{equation} \label{eq:13.1.5}
\rho y(b)+\delta y'(b)=k_{2}.
\end{equation}

\medskip
The assumptions  stated in this problem apply throughout this section and
won't be repeated. Note that we imposed conditions on $P_{0}$,
$P_{1}$, $P_{2}$, and $F$  on the \emph{closed} interval $[a,b]$, and we
are interested in solutions of \eqref{eq:13.1.3} on the closed interval.
This is different from the situation considered in Chapter~5,
where we imposed conditions on $P_{0}$, $P_{1}$, $P_{2}$, and $F$  on the
\emph{open} interval $(a,b)$ and we were interested in solutions on the
open interval. There
is really no problem here;   we can always extend  $P_{0}$, $P_{1}$,
$P_{2}$,
and $F$ to an open interval $(c,d)$ (for example, by defining them to be
constant on $(c,d]$ and $[b,d)$) so that they are continuous and $P_{0}$
has no zeros on $[c,d]$. Then we can apply the theorems from
Chapter~5 to the equation
$$
y''+\frac{P_{1}(x)}{P_{0}(x)}y'+
\frac{P_{2}(x)}{P_{0}(x)}y=\frac{F(x)}{P_{0}(x)}
$$
on $(c,d)$ to draw conclusions about solutions of \eqref{eq:13.1.3} on $[a,b]$.

We call $a$ and $b$ \emph{boundary points.}  The conditions \eqref{eq:13.1.4}
and \eqref{eq:13.1.5}  are \emph{boundary conditions}, and the problem is a
\emph{two-point boundary value problem} or, for simplicity, a
\emph{boundary value problem}. (We used  similar terminology in
Chapter~12 with a different meaning; both meanings are in common usage.)
  We require \eqref{eq:13.1.2} to insure that we're imposing a sensible
condition at each boundary point. For example, if $\alpha^{2}+\beta^{2}=0$
then $\alpha=\beta=0$, so $\alpha y(a)+\beta y'(a)=0$ for all choices of
$y(a)$ and $y'(a)$. Therefore  \eqref{eq:13.1.4} is an impossible condition if
$k_{1}\ne0$, or no condition at all if $k_{1}=0$.

We abbreviate \eqref{eq:13.1.1} as $Ly=F$, where
$$
Ly=P_{0}(x)y''+P_{1}(x)y'+P_{0}(x)y,
$$
and we denote
$$
B_{1}(y)=\alpha y(a)+\beta y'(a)
\text{\;    and\; \;  }
B_{2}(y)=\rho y(b)+\delta y'(b).
$$
We combine \eqref{eq:13.1.3}, \eqref{eq:13.1.4}, and \eqref{eq:13.1.5} as
\begin{equation} \label{eq:13.1.6}
Ly=F,\quad \quad B_{1}(y)=k_{1},\quad \quad B_{2}(y)=k_{2}.
\end{equation}
This boundary value problem is \emph{homogeneous} if $F=0$
and $k_{1}=k_{2}=0$;   otherwise it's \emph{nonhomogeneous}.

We leave it to you (Exercise~\ref{exer:13.1.1}) to verify that $B_{1}$
and $B_{2}$ are linear operators; that is, if $c_{1}$  and $c_{2}$
are constants then
\begin{equation} \label{eq:13.1.7}
B_{i}(c_{1}y_{1}+c_{2}y_{2})=c_{1}B_{i}(y_{1})+c_{2}B_{i}(y_{2}),\quad
i=1,2.
\end{equation}
The next three examples show that the question of existence and uniqueness
for solutions of boundary value problems is more complicated than
 for initial value problems.

\begin{example}\label{example:13.1.1}    \rm
Consider the boundary value problem
$$
y''+y=1,\quad y(0)=0, \quad y(\pi/2)=0.
$$
The general solution of $y''+y=1$ is
$$
y=1+c_{1}\sin x+c_{2} \cos x,
$$
so $y(0)=0$ if  and only if $c_{2}=-1$ and $y(\pi/2)=0$
if and only if $c_{1}=-1$. Therefore
$$
 y =1-\sin x -\cos x
$$
is the unique solution of the boundary value problem.
\end{example}

\begin{example}\label{example:13.1.2}    \rm
Consider the boundary value problem
$$
y''+y=1,\quad y(0)=0, \quad y(\pi)=0.
$$
Again, the general solution of $y''+y=1$ is
$$
y=1+c_{1}\sin x+c_{2} \cos x,
$$
so $y(0)=0$ if  and only if $c_{2}=-1$, but $y(\pi)=0$ if and only if
$c_{2}=1$. Therefore the boundary value problem has no solution.
\end{example}

\begin{example}\label{example:13.1.3}    \rm
Consider the boundary value problem
$$
y''+y=\sin 2x,\quad y(0)=0, \quad y(\pi)=0.
$$
You can use the method of undetermined coefficients
(Section~5.5) to find that the general solution of
$y''+y=\sin 2x$  is
$$
y=-\frac{\sin2x}{3}+c_{1}\sin x+c_{2}\cos x.
$$
The boundary conditions $y(0)=0$ and $y(\pi)=0$ both require that
$c_{2}=0$, but they don't restrict $c_{1}$. Therefore the boundary value
problem has infinitely many solutions
$$
y=-\frac{\sin2x}{3}+c_{1}\sin x,
$$
where $c_{1}$  is arbitrary.
\end{example}

\begin{theorem}\color{blue} \label{thmtype:13.1.1}
If $z_{1}$ and $z_{2}$ are solutions of $Ly=0$  such that either
$B_{1}(z_{1})=B_{1}(z_{2})=0$ or
$B_{2}(z_{1})=B_{2}(z_{2})=0$,  then $\{z_{1},z_{2}\}$  is linearly
dependent. Equivalently, if\,  $\{z_{1},z_{2}\}$ is linearly
independent, then
$$
B_{1}^{2}(z_{1})+B_{1}^{2}(z_{2})\ne 0 \text{\;    and\; \;  }
B_{2}^{2}(z_{1})+B_{2}^{2}(z_{2})\ne 0.
$$
\end{theorem}

\medskip
\proof
Recall that $B_{1}(z)=\alpha z(a)+\beta z'(a)$ and
$\alpha^{2}+\beta^{2}\ne0$.
Therefore, if $B_{1}(z_{1})=B_{1}(z_{2})=0$ then $(\alpha,\beta)$
is a nontrivial solution of  the system
\begin{eqnarray*}
\alpha z_{1}(a)+\beta z_{1}'(a)&=& 0\\
\alpha z_{2}(a)+\beta z_{}'(a)&=& 0.
\end{eqnarray*}
This  implies that
$$
z_{1}(a)z_{2}'(a)-z_{1}'(a)z_{2}(a)=0,
$$
so $\{z_{1},z_{2}\}$ is linearly dependent, by
Theorem~\ref{thmtype:5.1.6}. We leave it to you to show
that $\{z_{1},z_{2}\}$ is linearly dependent   if
$B_{2}(z_{1})=B_{2}(z_{2})=0$.

\begin{theorem}\color{blue} \label{thmtype:13.1.2}
The following statements are equivalent; that is, they are either
all true or all false.
\begin{alist}
\item    %(a)
 There's a fundamental set $\{z_{1},z_{2}\}$ of solutions of
$Ly=0$ such that
\begin{equation} \label{eq:13.1.8}
B_{1}(z_{1})B_{2}(z_{2})-B_{1}(z_{2})B_{2}(z_{1})\ne0.
\end{equation}
\item %(b)
If  $\{y_{1},y_{2}\}$ is a fundamental set of solutions of $Ly=0$  then
\begin{equation} \label{eq:13.1.9}
B_{1}(y_{1})B_{2}(y_{2})-B_{1}(y_{2})B_{2}(y_{1})\ne0.
\end{equation}
\item % (c)
For each continuous $F$ and pair of constants $(k_{1},k_{2})$,  the
boundary value problem
$$
Ly=F, \quad B_{1}(y)=k_{1},\quad B_{2}(y)=k_{2}
$$
has a unique solution.
\item % (d)
The homogeneous boundary value problem
\begin{equation} \label{eq:13.1.10}
Ly=0,\quad  B_{1}(y)=0, \quad B_{2}(y)=0
\end{equation}
has only the trivial solution $y=0$.
\item %(e)
The homogeneous equation $Ly=0$ has linearly independent solutions
$z_{1}$ and $z_{2}$ such that $B_{1}(z_{1})=0$ and $B_{2}(z_{2})=0$.
\end{alist}
\end{theorem}

\proof
We'll show  that
$$
\part{a}\implies
\part{b}\implies
\part{c}\implies
\part{d}\implies
\part{e}\implies
\part{a}.
$$

\part{a} $\implies$ \part{b}:
Since $\{z_{1},z_{2}\}$ is a fundamental set of solutions for $Ly=0$,
there are
constants
$a_{1}$, $a_{2}$, $b_{1}$, and $b_{2}$  such that
\begin{equation} \label{eq:13.1.11}
\begin{aligned}
y_{1} &=a_{1}z_{1}+a_{2}z_{2}\\
y_{2} &=b_{1}z_{1}+b_{2}z_{2}.
\end{aligned}
\end{equation}
Moreover,
\begin{equation} \label{eq:13.1.12}
\left|\begin{array}{ccccccc}
a_{1}&a_{2}\\b_{1}&b_{2}
\end{array}\right|\ne0.
\end{equation}
because if this determinant were zero, its rows would be linearly
dependent and therefore
$\{y_{1},y_{2}\}$ would be linearly dependent, contrary to our assumption
that  $\{y_{1},y_{2}\}$  is a fundamental set of solutions of $Ly=0$.
From \eqref{eq:13.1.7} and \eqref{eq:13.1.11},
$$
\left[\begin{array}{ccccccc}
B_{1}(y_{1})&B_{2}(y_{1})\\B_{1}(y_{2})&B_{2}(y_{2})
\end{array}\right]=
\left[\begin{array}{ccccccc}
a_{1}&a_{2}\\b_{1}&b_{2}
\end{array}\right]
\left[\begin{array}{ccccccc}
B_{1}(z_{1})&B_{2}(z_{1})\\
B_{1}(z_{2})&B_{2}(z_{2})
\end{array}\right].
$$
Since the determinant of a product of matrices is the product of the
determinants of the matrices, \eqref{eq:13.1.8} and \eqref{eq:13.1.12} imply
\eqref{eq:13.1.9}.

\part{b} $\implies$ \part{c}:
Since $\{y_{1},y_{2}\}$ is a fundamental set of solutions of
$Ly=0$, the general solution of  $Ly=F$ is
$$
y=y_{p}+c_{1}y_{1}+c_{2}y_{2},
$$
where $c_{1}$ and $c_{2}$ are arbitrary constants and  $y_{p}$ is a
particular solution of
$Ly=F$. To satisfy the boundary conditions, we must choose $c_{1}$ and
$c_{2}$ so that
\begin{eqnarray*}
k_{1}&=&B_{1}(y_{p})+c_{1}B_{1}(y_{1})+c_{2}B_{1}(y_{2})\\
k_{2}&=&B_{2}(y_{p})+c_{1}B_{2}(y_{1})+c_{2}B_{2}(y_{2}),
\end{eqnarray*}
(recall \eqref{eq:13.1.7}), which is equivalent to
$$
\begin{aligned}
c_{1}B_{1}(y_{1})+c_{2}B_{1}(y_{2}) &= k_{1}-B_{1}(y_{p}) \\
c_{1}B_{2}(y_{1})+c_{2}B_{2}(2_{2}) &= k_{2}-B_{2}(y_{p}).
\end{aligned}
$$
From  \eqref{eq:13.1.9}, this system always has a  unique solution
 $(c_{1},c_{2})$.

\part {c} $\implies$ \part{d}: Obviously, $y=0$ is a solution of
\eqref{eq:13.1.10}. From \part{c} with $F=0$ and $k_{1}=k_{2}=0$, it's the only
solution.

\part{d} $\implies$ \part{e}: Let  $\{y_{1},y_{2}\}$ be a fundamental
system for $Ly=0$ and  let
$$
z_{1}=B_{1}(y_{2})y_{1}-B_{1}(y_{1})y_{2}\text{\;    and\; \;  }
z_{2}=B_{2}(y_{2})y_{1}-B_{2}(y_{1})y_{2}.
$$
Then $B_{1}(z_{1})=0$  and $B_{2}(z_{2})=0$. To see
that $z_{1}$ and $z_{2}$ are linearly independent, note that
\begin{eqnarray*}
a_{1}z_{1}+a_{2}z_{2}&=&
a_{1}[B_{1}(y_{2})y_{1}-B_{1}(y_{1})y_{2}]+
a_{2}[B_{2}(y_{2})y_{1}-B_{2}(y_{1})y_{2}]\\
&=&
[B_{1}(y_{2})a_{1}+B_{2}(y_{2})a_{2}]y_{1}-[B_{1}(y_{1})a_{1}+B_{2}(y_{1})a_{2}]y_{2}.
\end{eqnarray*}
Therefore, since $y_{1}$ and $y_{2}$ are linearly independent,
$a_{1}z_{1}+a_{2}z_{2}=0$ if and only if
$$
\left[\begin{array}{ccccccc}
B_{1}(y_{1})&B_{2}(y_{1})\\
B_{1}(y_{2})&B_{2}(y_{2})
\end{array}\right]
\left[\begin{array}{ccccccc}
a_{1}\\a_{2}
\end{array}\right]=
\left[\begin{array}{ccccccc}
0\\0
\end{array}\right].
$$
If this system has a nontrivial solution then so does the system
$$
\left[\begin{array}{ccccccc}
B_{1}(y_{1})&B_{1}(y_{2})\\
B_{2}(y_{1})&B_{2}(y_{2})
\end{array}\right]
\left[\begin{array}{ccccccc}
c_{1} \\c_{2}
\end{array}\right]=
\left[\begin{array}{ccccccc}
0\\0
\end{array}\right].
$$
This and \eqref{eq:13.1.7} imply that $y=c_{1}z_{1}+c_{2}z_{2}$ is
a nontrivial solution of \eqref{eq:13.1.10}, which contradicts~\part{d}.

\part{e} $\implies$ \part{a}.
Theorem~\ref{thmtype:13.1.1} implies that if  $B_{1}(z_{1})=0$  and
$B_{2}(z_{2})=0$ then $B_{1}(z_{2})\ne0$ and $B_{2}(z_{1})\ne0$.
This implies~\eqref{eq:13.1.8}, which completes the proof.

\begin{example}\label{example:13.1.4} \rm
Solve the boundary value problem
\begin{equation} \label{eq:13.1.13}
x^{2}y''-2xy'+2y-2x^{3}=0,\quad y(1)=4,\quad y'(2)=3,
\end{equation}
given that $\{x,x^{2}\}$ is a fundamental set of solutions
of the complementary equation
\end{example}

\solution
Using variation of parameters (Section~5.7), you can show
that
$y_{p}=x^{3}$ is a solution of the complementary equation
$$
x^{2}y''-2xy'+2y=2x^{3}=0.
$$
Therefore the solution of
\eqref{eq:13.1.13} can be written as
$$
y=x^{3}+c_{1}x+c_{2}x^{2}.
$$
Then
$$
y'=3x^{2}+c_{1}+2c_{2}x,
$$
and imposing the boundary conditions yields the system
\begin{eqnarray*}
c_{1}+\phantom{4}c_{2}&=&\phantom{-}3\\ c_{1}+4c_{2}&=&-9,
\end{eqnarray*}
so $c_{1}=7$ and $c_{2}=-4$.
Therefore
$$
y=x^{3}+7x-4x^{2}
$$
is the unique solution of \eqref{eq:13.1.13}

\begin{example}\label{example:13.1.5} \rm
Solve the boundary value problem
$$
y''-7y'+12y=4e^{2x},\quad y(0)=3,\quad y(1)=5e^{2}.
$$
\end{example}

\solution
From Example~\ref{example:5.4.1}, $y_{p}=2e^{2x}$ is a particular solution
of
\begin{equation} \label{eq:13.1.14}
y''-7y'+12y=4e^{2x}.
\end{equation}
Since $\{e^{3x},e^{4x}\}$ is a fundamental set for the complementary
equation, we could write the solution of \eqref{eq:13.1.13} as
$$
y=2e^{2x}+c_{1}e^{3x}+c_{2}e^{4x}
$$
and determine $c_{1}$ and $c_{2}$ by imposing the boundary conditions.
However, this would lead to some tedious algebra, and the form of the
solution would be very unappealing. (Try it!) In this case it's convenient
to use the fundamental system $\{z_{1},z_{2}\}$ mentioned in
Theorem~\ref{thmtype:13.1.2}\part{e}; that is, we choose
$\{z_{1},z_{2}\}$ so that $B_{1}(z_{1})=z_{1}(0)=0$ and
$B_{2}(z_{2})=z_{2}(1)=0$. It is easy to see that
$$
z_{1}=e^{3x}-e^{4x}\text{\;    and\; \;  }z_{2}=e^{3(x-1)}-e^{4(x-1)}
$$
satisfy these requirements.
Now we write the solution of \eqref{eq:13.1.14} as
$$
y=2e^{2x}+c_{1}\left(e^{3x}-e^{4x}\right)+c_{2}\left(e^{3(x-1)}-e^{4(x-1)}\right).
$$
Imposing the boundary conditions $y(0)=3$  and $y(1)=5e^{2}$ yields
$$
3=2+c_{2}e^{-4}(e-1)\text{\;    and\; \;  }
5e^{2}=2e^{2}+c_{1}e^{3}(1-e).
$$
Therefore
$$
c_{1}=\frac{3}{e(1-e)},\quad c_{2}=\frac{e^{4}}{e-1},
$$
and
$$
y=2e^{2x}+\frac{3}{e(1-e)}(e^{3x}-e^{4x})+
\frac{e^{4}}{e-1}(e^{3(x-1)}-e^{4(x-1)}).
$$

Sometimes  it's useful to have a formula for the solution of a general
boundary problem. Our next theorem addresses this question.

\begin{theorem}\color{blue} \label{thmtype:13.1.3}
Suppose the homogeneous boundary value problem
\begin{equation} \label{eq:13.1.15}
Ly=0,\quad B_{1}(y)=0,\quad B_{2}(y)=0
\end{equation}
has only the trivial solution. Let $y_{1}$ and $y_{2}$ be linearly
independent solutions   of $Ly=0$ such that $B_{1}(y_{1})=0$ and
$B_{2}(y_{2})=0$, and let
$$
W=y_{1}y_{2}'-y_{1}'y_{2}.
$$
Then the unique solution of
\begin{equation} \label{eq:13.1.16}
Ly=F,\quad B_{1}(y)=0,\quad B_{2}(y)=0
\end{equation}
is
\begin{equation} \label{eq:13.1.17}
y(x)= y_{1}(x)\int_{x}^{b}\frac{F(t)y_{2}(t)}{P_{0}(t)W(t)}\,dt+
y_{2}(x)\int_{a}^{x}\frac{F(t)y_{1}(t)}{P_{0}(t)W(t)}\,dt.
\end{equation}
\end{theorem}

\medskip
\proof In Section~5.7  we saw that if
\begin{equation} \label{eq:13.1.18}
y=u_{1}y_{1}+u_{2}y_{2}
\end{equation}
where
\begin{eqnarray*}
u_{1}'y_{1}+u_{2}'y_{2}&=&0\\
u_{1}'y_{1}'+u_{2}'y_{2}'&=&F,
\end{eqnarray*}
then  $Ly=F$. Solving for $u_{1}'$ and $u_{2}'$ yields
$$
u_{1}'=-\frac{Fy_{2}}{P_{0}W}\text{\;    and\; \;  }
u_{2}'=\frac{Fy_{1}}{P_{0}W},
$$
which hold if
$$
u_{1}(x)=\int_{x}^{b}\frac{F(t)y_{2}(t)}{P_{0}(t)W(t)}\,dt
\text{\;    and\; \;  }
u_{2}(x)=\int_{a}^{x}\frac{F(t)y_{1}(t)}{P_{0}(t)W(t)}\,dt.
$$
This and \eqref{eq:13.1.18} show that \eqref{eq:13.1.17} is a solution of $Ly=F$.
Differentiating \eqref{eq:13.1.17}  yields
\begin{equation} \label{eq:13.1.19}
y'(x)=y_{1}'(x)\int_{x}^{b}\frac{F(t)y_{2}(t)}{P_{0}(t)W(t)}\,dt+
y_{2}'(x)\int_{a}^{x}\frac{F(t)y_{1}(t)}{P_{0}(t)W(t)}\,dt.
\end{equation}
(Verify.) From \eqref{eq:13.1.17} and \eqref{eq:13.1.19},
$$
B_{1}(y)=B_{1}(y_{1})\int_{a}^{b}\frac{F(t)y_{2}(t)}{P_{0}(t)W(t)}\,dt=0
$$
because $B_{1}(y_{1})=0$, and
$$
B_{2}(y)=B_{2}(y_{2})\int_{a}^{b}\frac{F(t)y_{1}(t)}{P_{0}(t)W(t)}\,dt=0
$$
because $B_{2}(y_{2})=0$. Hence, $y$  satisfies \eqref{eq:13.1.16}.
This completes the proof.

We  can rewrite \eqref{eq:13.1.17} as
\begin{equation} \label{eq:13.1.20}
y=\int_{a}^{b} G(x,t)F(t)\,dt,
\end{equation}
where
$$
G(x,t)=
\begin{cases}
\dst{\frac{y_{1}(t)y_{2}(x)}{P_{0}(t)W(t)}}&,\quad a\le t\le x,\\[.1in]
\dst{\frac{y_{1}(x)y_{2}(t)}{P_{0}(t)W(t)}} &,\quad x\le t\le b.
\end{cases}.
$$

This is the   \emph{Green's function} for the boundary value problem
\eqref{eq:13.1.16}.
The Green's function  is related to the boundary value problem
\eqref{eq:13.1.16} in much the
same way that the inverse of a square matrix ${\mathbf A}$ is related to
the linear algebraic system ${\mathbf y}={\mathbf A}{\mathbf x}$;   just as
we substitute the given vector ${\mathbf y}$ into the formula
${\mathbf x}={\mathbf A}^{-1}{\mathbf y}$ to solve
 ${\mathbf y} ={\mathbf Ax}$, we substitute the given
function $F$ into the formula \eqref{eq:13.1.20} to obtain the solution of
\eqref{eq:13.1.16}. The analogy goes further:
just as ${\mathbf A}^{-1}$ exists if and only if
 ${\mathbf A}{\mathbf x} ={\mathbf 0}$
has only the trivial solution, the boundary value problem  \eqref{eq:13.1.16}
has a Green's function if and only the homogeneous boundary value problem
\eqref{eq:13.1.15}  has only the trivial solution.

We leave it to you (Exercise~\ref{exer:13.1.32}) to show that the
assumptions of Theorem~\ref{thmtype:13.1.3} imply that the unique solution
of the boundary value problem
$$
Ly=F,\quad B_{1}(y)=k_{1},\quad B_{2}(y)=k_{2}
$$
is
$$
y(x)=\int_{a}^{b}G(x,t)F(t)\,dt +\frac{k_{2}}{B_{2}(y_{1})}y_{1}+
\frac{k_{1}}{B_{1}(y_{2})}y_{2}.
$$

\begin{example}\label{example:13.1.6}\rm
Solve the boundary value problem
\begin{equation} \label{eq:13.1.21}
y''+y=F(x).\quad y(0)+y'(0)=0,\quad y(\pi)-y'(\pi)=0,
\end{equation}
and find the Green's function for this problem.
\end{example}

\solution
Here
$$
B_{1}(y)=y(0)+y'(0)   \text{\;    and\; \;  } B_{2}(y)=y(\pi)-y'(\pi).
$$
Let $\{z_{1},z_{2}\}=\{\cos x,\sin x\}$, which is a fundamental set of
solutions of $y''+y=0$. Then
\begin{eqnarray*}
B_{1}(z_{1})&=&(\cos x-\sin x)\big|_{x=0}=\phantom{-}1\\
B_{2}(z_{1})&=&(\cos x+\sin x)\big|_{x=\pi}=-1
\end{eqnarray*}
and
\begin{eqnarray*}
B_{1}(z_{2})&=&(\sin x+\cos x)\big|_{x=0}=1\\
B_{2}(z_{2})&=&(\sin x-\cos x)\big|_{x=\pi}=1.
\end{eqnarray*}
Therefore
$$
B_{1}(z_{1})B_{2}(z_{2})-B_{1}(z_{2})B_{2}(z_{1})=2,
$$
so Theorem~\ref{thmtype:13.1.2} implies that \eqref{eq:13.1.21} has a unique
solution. Let
$$
y_{1}=B_{1}(z_{2})z_{1}-B_{1}(z_{1})z_{2}= \cos x-\sin x
$$
 and
$$
y_{2}=B_{2}(z_{2})z_{1}-B_{2}(z_{1})z_{2}=  \cos x+\sin x.
$$
Then $B_{1}(y_{1})=0$, $B_{2}(y_{2})=0$, and the Wronskian of
$\{y_{1},y_{2}\}$ is
$$
W(x)=
\left|\begin{array}{rrccccc}
\cos x-\sin x& \cos x+ \sin x\\
-\sin x-\cos x&-\sin x+\cos x
\end{array}\right|=2.
$$
Since $P_{0}=1$, \eqref{eq:13.1.17} yields the solution
\begin{eqnarray*}
y(x)&=&\dst{\frac{\cos x-\sin x}{2}}\int_{x}^{\pi}F(t)(\cos t+\sin t)\,dt\\
&&+\dst{\frac{\cos x +\sin x}{2}}\int_{0}^{x}F(t)(\cos t-\sin t)\,dt.
\end{eqnarray*}
The Green's function is
$$
G(x,t)=
\begin{cases}
\dst{\frac{(\cos t-\sin t)(\cos x+\sin x)}{2}},&\quad  0\le t\le
x,\\[.1in]
\dst{\frac{(\cos x-\sin x)(\cos t+\sin t)}{2}},&\quad x\le t\le \pi.
\end{cases}
$$

We'll now consider the situation not covered by Theorem~\ref{thmtype:13.1.3}.

\begin{theorem}\color{blue} \label{thmtype:13.1.4}
Suppose the homogeneous boundary value problem
\begin{equation} \label{eq:13.1.22}
Ly=0,\quad B_{1}(y)=0,\quad B_{2}(y)=0
\end{equation}
has a nontrivial solution $y_{1}$,  and let $y_{2}$ be  any solution of
$Ly=0$ that isn't  a constant multiple of $y_{1}$.  Let
$W=y_{1}y_{2}'-y_{1}'y_{2}$. If
\begin{equation} \label{eq:13.1.23}
\int_{a}^{b}\frac{F(t)y_{1}(t)}{P_{0}(t)W(t)}\,dt=0,
\end{equation}
then the homogeneous boundary value problem
\begin{equation} \label{eq:13.1.24}
Ly=F, \quad B_{1}(y)=0,\quad B_{2}(y)=0
\end{equation}
has infinitely many solutions, all of the form $y=y_{p}+c_{1}y_{1}$,
where
$$
y_{p}=y_{1}(x)
\int_{x}^{b}\frac{F(t)y_{2}(t)}{P_{0}(t)W(t)}\,dt+
y_{2}(x)
\int_{a}^{x}\frac{F(t)y_{1}(t)}{P_{0}(t)W(t)}\,dt
$$
and $c_{1}$ is a constant. If
$$
\int_{a}^{b}\frac{F(t)y_{1}(t)}{P_{0}(t)W(t)}\,dt \ne0,
$$
then \eqref{eq:13.1.24}  has no solution.
\end{theorem}

\proof
From the proof of Theorem~\ref{thmtype:13.1.3}, $y_{p}$ is a particular solution
of $Ly=F$,
and
$$
y_{p}'(x)=y_{1}'(x)\int_{x}^{b}\frac{F(t)y_{2}(t)}{P_{0}(t)W(t)}\,dt
+y_{2}'(x)\int_{a}^{x}\frac{F(t)y_{1}(t)}{P_{0}(t)W(t)}\,dt.
$$
Therefore the general solution of \eqref{eq:13.1.22} is of the form
$$
y=y_{p}+c_{1}y_{1}+c_{2}y_{2},
$$
where $c_{1}$ and $c_{2}$  are constants.  Then
\begin{eqnarray*}
B_{1}(y)&=&
B_{1}(y_{p}+c_{1}y_{1}+c_{2}y_{2})=
B_{1}(y_{p})+c_{1}B_{1}(y_{1})+c_{2}B_{1}y_{2}\\
&=&B_{1}(y_{1})\int_{a}^{b}\frac{F(t)y_{2}(t)}{P_{0}(t)W(t)}\,dt
+c_{1}B_{1}(y_{1})+c_{2}B_{1}(y_{2})\\
&=&c_{2}B_{1}(y_{2})
\end{eqnarray*}
Since $B_{1}(y_{1})=0$, Theorem~\ref{thmtype:13.1.1} implies that
$B_{1}(y_{2})\ne0$; hence, $B_{1}(y)=0$ if and only if $c_{2}=0$.
Therefore $y=y_{p}+c_{1}y_{1}$ and
\begin{eqnarray*}
B_{2}(y)&=&B_{2}(y_{p}+c_{1}y_{1})
=B_{2}(y_{2})\int_{a}^{b}\frac{F(t)y_{1}(t)}{P_{0}(t)W(t)}\,dt+c_{1}B_{2}(y_{1})\\
&=&B_{2}(y_{2})\int_{a}^{b}\frac{F(t)y_{1}(t)}{P_{0}(t)W(t)}\,dt,
\end{eqnarray*}
since $B_{2}(y_{1})=0$.  From Theorem~\ref{thmtype:13.1.1},
$B_{2}(y_{2})\ne0$ (since $B_{2}(y_{1}=0$).
Therefore $Ly=0$   if and only if  \eqref{eq:13.1.23} holds.
This completes the proof.

\begin{example}\label{example:13.1.7} \rm
Applying Theorem~\ref{thmtype:13.1.4} to the boundary value problem
\begin{equation} \label{eq:13.1.25}
y''+y=F(x),\quad y(0)=0,\quad y(\pi)=0
\end{equation}
explains the  Examples~\ref{example:13.1.2} and \ref{example:13.1.3}. The
complementary equation $y''+y=0$ has the linear independent solutions
$y_{1}=\sin x$ and $y_{2}=\cos x$, and $y_{1}$ satisfies both boundary
conditions. Since $P_{0}=1$ and
$$
W=
\left|\begin{array}{crccccc}
\sin x&\cos x \\\cos x&-\sin x
\end{array}\right|=-1,
$$
\eqref{eq:13.1.23} reduces to
$$
\int_{0}^{\pi} F(x)\sin x\,dx=0.
$$
From Example~\ref{example:13.1.2}, $F(x)=1$ and
$$
\int_{0}^{\pi} F(x)\sin x\,dx=\int_{0}^{\pi} \sin x\,dx=2,
$$
so      Theorem~\ref{thmtype:13.1.3}  implies that \eqref{eq:13.1.25} has no
solution.  In Example~\ref{example:13.1.3},
$$
F(x)=\sin 2x=2\sin x \cos x
$$
and
$$
\int_{0}^{\pi}F(x)\sin x\,dx=2\int_{0}^{\pi}\sin^{2}x\cos x\,dx
=\frac{2}{3}\sin^{3}x\bigg|_{0}^{\pi}=0,
$$
so Theorem~\ref{thmtype:13.1.3} implies that \eqref{eq:13.1.25} has infinitely
many solutions, differing by constant multiples of $y_{1}(x)=\sin x$.
\end{example}

\exercises

\begin{exerciselist}

\item\label{exer:13.1.1}
Verify that $B_{1}$ and $B_{2}$ are linear operators; that is, if $c_{1}$
and $c_{2}$  are constants then
$$
B_{i}(c_{1}y_{1}+c_{2}y_{2})=c_{1}B_{i}(y_{1})+c_{2}B_{i}(y_{2}),\quad
i=1,2.
$$

\exercisetext{In Exercises~\ref{exer:13.1.2}--\ref{exer:13.1.7} solve the boundary
value problem.}

\item\label{exer:13.1.2} $y''-y=x$, \quad  $y(0)=-2$,\quad  $y(1)=1$

\item\label{exer:13.1.3} $y''=2-3x$, \quad  $y(0)=0$,\quad  $y(1)-y'(1)=0$

\item\label{exer:13.1.4} $y''-y=x$, \quad  $y(0)+y'(0)=3$,\quad  $y(1)-y'(1)=2$

\item\label{exer:13.1.5} $y''+4y=1$, \quad  $y(0)=3$,\quad
 $y(\pi/2)+y'(\pi/2)=-7$

\item\label{exer:13.1.6} $y''-2y'+y=2e^{x}$, \quad  $y(0)-2y'(0)=3$,\quad
 $y(1)+y'(1)=6e$

\item\label{exer:13.1.7} $y''-7y'+12y=4e^{2x}$, \quad  $y(0)+y'(0)=8$,\quad
 $y(1)=-7e^{2}$ (see Example~\ref{example:13.1.5})

\item\label{exer:13.1.8} State a condition on $F$ such that the boundary
value problem
$$
y''=F(x), \quad  y(0)=0, \quad y(1)-y'(1)=0
$$
has a  solution, and find all solutions.

\item\label{exer:13.1.9}
\begin{alist}
\item % (a)
State  a condition on $a$ and $b$ such that the boundary
value problem
\begin{equation}
y''+y=F(x),\quad y(a)=0,\quad y(b)=0     \tag{A}
\end{equation}
has a unique solution for every continuous $F$, and find the solution
by the method used to prove Theorem~\ref{thmtype:13.1.3}
\item % (b)
In the case where $a$ and $b$ don't satisfy the condition you gave for
\part{a}, state necessary and sufficient     on $F$ such that (A) has
a solution, and find all solutions by the method used to prove
Theorem~\ref{thmtype:13.1.4}.
\end{alist}

\item\label{exer:13.1.10}
Follow the instructions in Exercise~\ref{exer:13.1.9} for the boundary value
problem
$$
y''+y=F(x),\quad y(a)=0,\quad y'(b)=0.
$$

\item\label{exer:13.1.11}
Follow the instructions in Exercise~\ref{exer:13.1.9} for the boundary value
problem
$$
y''+y=F(x),\quad y'(a)=0,\quad y'(b)=0.
$$

\exercisetext{In Exercises~\ref{exer:13.1.12}--\ref{exer:13.1.15} find a formula for
the solution of the  boundary problem by the method used to prove
Theorem~\ref{thmtype:13.1.3}. Assume that $a<b$.}

\item\label{exer:13.1.12}
$y''-y=F(x)$,\quad  $y(a)=0$,\quad  $y(b)=0$

\item\label{exer:13.1.13}
$y''-y=F(x)$,\quad  $y(a)=0$,\quad  $y'(b)=0$

\item\label{exer:13.1.14}
$y''-y=F(x)$,\quad  $y'(a)=0$,\quad  $y'(b)=0$

\item\label{exer:13.1.15}
$y''-y=F(x)$,\quad  $y(a)-y'(a)=0$,\quad  $y(b)+y'(b)=0$

\exercisetext{In Exercises~\ref{exer:13.1.16}--\ref{exer:13.1.19} find all values of
$\omega$ such that
 boundary problem has a unique solution, and find the solution  by the
method used to prove Theorem~\ref{thmtype:13.1.3}. For other values of $\omega$,
find conditions on $F$ such that the problem has a solution,
 and find all solutions  by the
method used to prove Theorem~\ref{thmtype:13.1.4}.}

\item\label{exer:13.1.16}
$y''+ \omega^{2}y=F(x)$,\quad $y(0)=0$,\quad $y(\pi)=0$

\item\label{exer:13.1.17}
$y''+ \omega^{2}y=F(x)$,\quad $y(0)=0$,\quad $y'(\pi)=0$

\item\label{exer:13.1.18}
$y''+ \omega^{2}y=F(x)$,\quad $y'(0)=0$,\quad $y(\pi)=0$

\item\label{exer:13.1.19}
$y''+ \omega^{2}y=F(x)$,\quad $y'(0)=0$,\quad $y'(\pi)=0$

\item\label{exer:13.1.20}
Let $\{z_{1},z_{2}\}$ be a fundamental set of solutions of $Ly=0$.
Given that the homogeneous boundary value problem
$$
Ly=0, \quad B_{1}(y)=0,\quad B_{2}(y)=0
$$
has a nontrivial solution, express it explicitly in terms of $z_{1}$
and $z_{2}$.

\item\label{exer:13.1.21}
If the boundary value problem has a solution for every continuous  $F$,
then find  the Green's function for the problem and use it to write
an explicit formula for the solution. Otherwise, if the
boundary value problem does not have a solution for every continuous $F$,
find a necessary and sufficient condition on $F$ for the problem to have
a solution, and find all solutions. Assume that $a<b$.

\begin{alist}
\item % (a)
$y''=F(x)$, \quad $y(a)=0$,\quad $y(b)=0$

\item % (b)
$y''=F(x)$, \quad $y(a)=0$,\quad $y'(b)=0$

\item % (c)
$y''=F(x)$, \quad $y'(a)=0$,\quad $y(b)=0$

\item % (d)
$y''=F(x)$, \quad $y'(a)=0$,\quad $y'(b)=0$
\end{alist}

\item\label{exer:13.1.22}
Find the Green's function  for the boundary value problem
\begin{equation}
y''=F(x), \quad y(0)-2y'(0)=0, \quad y(1)+2y'(1)=0.   \tag{A}
\end{equation}
Then  use the Green's function to solve (A) with
\part{a} $F(x)=1$, \part{b} $F(x)=x$, and \part{c} $F(x)=x^{2}$.

\item\label{exer:13.1.23}
Find the Green's function  for the boundary value problem
\begin{equation}
x^{2}y''+xy'+(x^{2}-1/4)y=F(x), \quad y(\pi/2)=0,\quad y(\pi)=0,   \tag{A}
\end{equation}
given that
$$
y_{1}(x)=\frac{\cos x}{\sqrt{x}}
\text{\;    and\; \;  }
y_{2}(x)=\frac{\sin x}{\sqrt{x}}
$$
are solutions of the complementary equation. Then use the Green's function
to solve  (A) with \part{a} $F(x)=x^{3/2}$  and \part{b} $F(x)=x^{5/2}$.

\item\label{exer:13.1.24}
Find the Green's function  for the boundary value problem
\begin{equation}
x^{2}y''-2xy'+2y=F(x), \quad y(1)=0,\quad y(2)=0,   \tag{A}
\end{equation}
given that  $\{x,x^{2}\}$ is a fundamental set of solutions
of the complementary equation.
Then use the Green's function
to solve  (A) with \part{a} $F(x)=2x^{3}$  and \part{b} $F(x)~=~6x^{4}$.

\item\label{exer:13.1.25}
Find the Green's function  for the boundary value problem
\begin{equation}
x^{2}y''+xy'-y=F(x), \quad y(1)-2y'(1)=0,\quad y'(2)=0,   \tag{A}
\end{equation}
given that  $\{x,1/x\}$ is a fundamental set of solutions
of the complementary equation.
Then use the Green's function
to solve  (A) with \part{a} $F(x)=1$,   \part{b} $F(x)=x^{2}$, and
\part{c} $F(x)=x^{3}$.

\exercisetext{In Exercises~\ref{exer:13.1.26}--\ref{exer:13.1.30} find
necessary and sufficient conditions on $\alpha$, $\beta$, $\rho$,
and $\delta$ for the boundary value
 problem to have a  unique solution for every continuous $F$, and
find the Green's function.}

\item\label{exer:13.1.26}
$y''=F(x)$, \quad $\alpha y(0)+\beta y'(0)=0$, \quad
$\rho y(1)+\delta y'(1)=0$

\item\label{exer:13.1.27}
$y''+y=F(x)$, \quad $\alpha y(0)+\beta y'(0)=0$, \quad
$\rho y(\pi)+\delta y'(\pi)=0$

\item\label{exer:13.1.28}
$y''+y=F(x)$, \quad $\alpha y(0)+\beta y'(0)=0$, \quad
$\rho y(\pi/2)+\delta y'(\pi/2)=0$

\item\label{exer:13.1.29}
$y''-2y'+2y=F(x)$, \quad $\alpha y(0)+\beta y'(0)=0$, \quad
$\rho y(\pi)+\delta y'(\pi)=0$

\item\label{exer:13.1.30}
$y''-2y'+2y=F(x)$, \quad $\alpha y(0)+\beta y'(0)=0$, \quad
$\rho y(\pi/2)+\delta y'(\pi/2)=0$

\item\label{exer:13.1.31}
Find
necessary and sufficient conditions on $\alpha$, $\beta$, $\rho$,
and $\delta$ for the boundary value
 problem
\begin{equation}
y''-y=F(x), \quad \alpha y(a)+\beta y'(a)=0, \quad
\rho y(b)+\delta y'(b)=0   \tag{A}
\end{equation}
to have a  unique solution for every continuous $F$, and find the Green's
function for (A). Assume that $a<b$.

\item\label{exer:13.1.32}
Show that  the assumptions of Theorem~\ref{thmtype:13.1.3} imply that the unique
solution of
$$
Ly=F, \quad  B_{1}(y)=k_{1},\quad B_{2}(y)=f_{2}
$$
is
$$
y=\int_{a}^{b} G(x,t)F(t)\,dt +\frac{k_{2}}{B_{2}}(y_{1})y_{1}
+\frac{k_{1}}{B_{1}(y_{2})}y_{2}.
$$
\end{exerciselist}
